\section{降维与升维的对偶：在合适的维度做合适的事}
\label{sec:13-Machine-Learning/04-Dimensionality-Reduction/05}

你在给一个检索服务选方案。一亿条 768 维的向量，全量放内存装不下，必须压。一派主张 PCA 到 64 维，理由是数据本来就低秩；另一派主张随机投影到 128 维，理由是快、不用训练、新数据来了直接乘。第二个方案听起来不太严肃：一个与数据毫无关系的高斯矩阵，凭什么保证压完还能检索？更让人困惑的是同一个系统的上游——那些 768 维向量本身，是模型把几十维的原始特征一路升上去得到的。一边拼命往上升，一边拼命往下压，这个系统到底在干什么？

\subsection{随机投影几乎不挑方向，它只承诺距离}

先回答第二个方案。把每个 \(x\in\R^d\) 映射成 \(\Pi x/\sqrt k\)，其中 \(\Pi\in\R^{k\times d}\) 的元素独立取自标准正态。Johnson--Lindenstrauss 引理说：给定 \(n\) 个点和 \(\varepsilon\in(0,1)\)，只要

\[
k = O\bigl(\varepsilon^{-2}\log n\bigr),
\]

就能以高概率让所有点对的距离同时满足

\[
(1-\varepsilon)\norm{x_i-x_j}_2 \le \norm{\tfrac{1}{\sqrt k}\Pi(x_i-x_j)}_2 \le (1+\varepsilon)\norm{x_i-x_j}_2 .
\]

请注意 \(k\) 里没有 \(d\)。原始维度是 768 还是 76800，需要的目标维度一样，只取决于点数和你要的精度。证明的骨架不长：对固定的一个差向量，\(\norm{\Pi u}_2^2/k\) 是 \(k\) 个独立卡方项的平均，它的偏离概率随 \(k\) 指数衰减；再对 \(\binom{n}{2}\) 个点对做并集界，指数衰减压得住多项式个事件，取 \(k\) 与 \(\log n\) 成正比即可。条件对账：需要投影矩阵的分布是各向同性的（高斯、\(\pm1\) 随机符号、稀疏化的变体都行），需要事先固定点集——它保护的是这 \(n\) 个点，不是整个空间。

这件事的意义在于它给出了一种全新的``声明''。本单元开头说过，任何降维都必须声明哪些方向的信息可以丢；PCA 声明的是方差小的方向可以丢，所以它必须先看数据、算协方差、挑方向。随机投影不挑方向，它甚至不看数据一眼，它声明的是另一件事：\textbf{点对之间的距离要保住，别的一概不管}。

\begin{center}
\begin{tikzpicture}[>=stealth,scale=1.0]
% 面板一：随机投影
\draw[gray!60,->] (0,0) -- (3.3,0);
\draw[gray!60,->] (0,0) -- (0,3.3);
\draw[black!45] (0,0) -- (3.1,3.1);
\draw[dashed,black!35] (0,0) -- (3.1,3.57);
\draw[dashed,black!35] (0,0) -- (3.1,2.64);
\foreach \a/\b in {0.4/0.43, 0.55/0.52, 0.7/0.66, 0.9/0.95, 1.05/1.1,
                   1.2/1.13, 1.4/1.48, 1.55/1.5, 1.7/1.62, 1.9/2.0,
                   2.05/2.15, 2.2/2.12, 2.4/2.5, 2.7/2.6, 2.85/2.9}
  \fill[black!72] (\a,\b) circle (1.6pt);
\node[font=\scriptsize,anchor=north] at (1.65,-0.18) {原始距离};
\node[font=\scriptsize,rotate=90,anchor=south] at (-0.3,1.65) {投影后距离};
\node[font=\small] at (1.55,-0.75) {随机投影};
% 面板二：PCA
\begin{scope}[shift={(5.0,0)}]
\draw[gray!60,->] (0,0) -- (3.3,0);
\draw[gray!60,->] (0,0) -- (0,3.3);
\draw[black!45] (0,0) -- (3.1,3.1);
\foreach \a/\b in {0.4/0.15, 0.55/0.5, 0.7/0.62, 0.9/0.3, 1.05/0.35,
                   1.2/1.15, 1.4/0.5, 1.55/1.5, 1.7/1.6, 1.9/0.8,
                   2.05/0.9, 2.2/2.1, 2.4/1.2, 2.7/2.6, 2.85/1.4}
  \fill[black!72] (\a,\b) circle (1.6pt);
\node[font=\scriptsize,anchor=north] at (1.65,-0.18) {原始距离};
\node[font=\small] at (1.55,-0.75) {PCA 投影};
\end{scope}
\end{tikzpicture}
\end{center}

\emph{每个点是一对样本。随机投影把所有点对压在对角线附近的窄带里；PCA 让沿主方向分开的点对留在对角线上，让沿被丢方向分开的点对整体塌下来。}

图里两幅散点的差别就是两种声明的差别。右边那些远离对角线的点，恰恰是 PCA 主动放弃的——它们在被截掉的方向上分开，投影之后就重合了。这不是 PCA 的缺陷，是它的定义。反过来，左边那条窄带也不免费：随机投影对每个点对一视同仁，它没有借数据的低秩结构占任何便宜，所以在同样的精度要求下需要的维数往往比 PCA 多，而且它不给你可解释的坐标轴。

顺带说清它保什么、不保什么。距离保住了，内积和角度也就在同样的 \(\varepsilon\) 量级上受控，因为内积可由三条边长表出。但面积、体积、行列式这类多点的联合几何量没有类似保证，误差会随涉及的点数累积；矩阵的秩和谱结构更是不保——随机投影之后，原本的低秩结构一般不再是低秩。要用哪条性质，就得单独验哪条性质。

\subsection{反过来，升维让纠缠的结构散开}

现在看上游那一半。取两个同心圆，内圈一类外圈一类，二维平面上没有任何直线能把它们分开。补一维，

\[
(x_1,x_2)\ \longmapsto\ (x_1,\ x_2,\ x_1^2+x_2^2),
\]

两类立刻被一张水平面分开。数据一个比特都没多，多的只是坐标；但在新坐标下，一个线性模型就够了。Cover 在 1965 年把这件事写成定理：给定一组一般位置的点，随机二分中线性可分的比例随维数上升而上升。升维不创造信息，它把已有的信息摆成模型能直接使用的形式。

核方法把这条路走到极致：不显式构造 \(\phi(x)\)，只用核函数

\[
K(x,x') = \ip{\phi(x)}{\phi(x')}
\]

在高维（乃至无穷维）里做内积。SVM、核岭回归、Gaussian process 都在这个框架下工作，计算发生在高维，存储却只在样本数的规模上。深度网络遵循同样的节奏：中间层把通道数往上抬，在高维里做精细的交互与混合；输出层再压回类别数那么低，做决策。一升一降不是自相矛盾，是分工。

\subsection{高维只在有结构时才划算}

升维不是无条件的好事。当维数远超样本数而数据没有任何结构时，样本变得稀疏，距离趋于均等，任何依赖``谁离谁近''的方法都失去分辨力，计算成本还照单全收。上一节的流形假设、稀疏假设、低秩假设，正是让高维仍然可算的那些结构；数据是白噪声或加密串时，这些假设一条不成立，升维只剩坏处。

同一个提醒在诊断层面也成立。样本协方差的谱在 \(p/n\) 不小时会系统性地散开，看起来像有信号；判断的基准应当是 Marchenko--Pastur 上界而不是肉眼，见\hyperref[ch:094]{``随机矩阵理论：高维协方差的谱失真''}一章。这条基准对本节尤其相关：你在 768 维上算出的``前 64 个主成分解释了大部分方差''，有多少是结构、有多少是维数比造成的，需要一个零模型来对照。

\subsection{计算放在高维，理解放在低维}

把两个方向摆在一起，标准就清楚了。需要精细计算——梯度、内积、线性可分性——就往高维去，代价由结构假设兜底；需要人看懂形态——聚类、异常、决策边界——就往低维压，代价是不可逆的信息损失。同一条流水线上两件事可以同时发生，检索服务就是例子：上游模型在高维里做表示学习，下游索引在低维里做距离比较。

最后回到那个选型问题。真正的答案不在两个方案本身，而在评估方式：把 PCA 或随机投影当作流程的一部分，在下游指标上比较。降维一旦接进系统，问题就已经换了——模型面对的不再是原始特征，而是一组按某种准则挑出来的坐标；被丢掉的方向如果恰好携带你关心的信息，下游再强也补不回来。因此投影矩阵、均值、方差必须只用训练折估计，评估必须落在召回率、命中率这类端到端指标上，而不是重构误差。这与\hyperref[sec:11-Mathematical-Statistics/06-Resampling-and-Model-Selection/03]{``交叉验证评估的是完整训练流程''}一节说的是同一条纪律。

到这里，本单元讲的降维都是确定性的几何操作：给一个矩阵，输出一组坐标。它们回答不了另一类问题——这个低维表示有多少不确定性？缺失值怎么办？怎么给新样本一个概率？下一节把低秩结构写成生成模型，这些问题才有位置安放。
